\documentclass[12pt,a4paper]{report}
\usepackage[utf8]{inputenc}
\usepackage[T5]{fontenc}
\usepackage[vietnamese]{babel}
\usepackage[a4paper,top=2.5cm,bottom=2.5cm,left=3cm,right=2cm]{geometry}
\usepackage{amsmath,amssymb,mathtools}
\usepackage{booktabs}
\usepackage{caption}
\usepackage{graphicx}
\usepackage{tikz}
\usetikzlibrary{arrows.meta,positioning,calc,shapes.geometric}
\usepackage{setspace}
\usepackage{array}
\usepackage{makecell}
\usepackage[numbers]{natbib}
\usepackage[hidelinks]{hyperref}
\usepackage[capitalize,nameinlink]{cleveref}

\setstretch{1.25}
\setlength{\parindent}{1.25cm}
\setlength{\parskip}{0.15em}

% -----------------------------------------------------------------
% Thông tin bìa
% -----------------------------------------------------------------
\newcommand{\UniversityName}{ĐẠI HỌC QUỐC GIA HÀ NỘI}
\newcommand{\SchoolName}{TRƯỜNG ĐẠI HỌC CÔNG NGHỆ}
\newcommand{\ThesisTitle}{TĂNG CƯỜNG KHẢ NĂNG HIỂU Ý ĐỊNH NGƯỜI DÙNG\\ CHO ROBOT CỘNG TÁC DỰA TRÊN MÔ HÌNH VLA}
\newcommand{\ThesisTitleEN}{ENHANCING USER INTENTION UNDERSTANDING FOR COLLABORATIVE ROBOTS BASED ON VLA MODELS}
\newcommand{\StudentName}{........................................................}
\newcommand{\StudentID}{........................................................}
\newcommand{\SupervisorName}{........................................................}
\newcommand{\MajorName}{........................................................}
\newcommand{\ThesisYear}{2026}

\begin{document}

% ============================================================
% Trang bìa
% ============================================================
\begin{titlepage}
    \centering
    \vspace*{0.8cm}
    {\bfseries\large \UniversityName\par}
    \vspace{0.25cm}
    {\bfseries\large \SchoolName\par}

    \vspace{2cm}
    {\bfseries\Large ĐỒ ÁN TỐT NGHIỆP\par}

    \vspace{1.5cm}
    {\bfseries\LARGE \ThesisTitle\par}
    \vspace{0.5cm}
    {\large \ThesisTitleEN\par}

    \vfill
    \begin{tabular}{>{\bfseries}p{4.8cm}p{8.3cm}}
        Sinh viên: & \StudentName \\[0.35cm]
        Mã sinh viên: & \StudentID \\[0.35cm]
        Ngành/Chương trình: & \MajorName \\[0.35cm]
        Giảng viên hướng dẫn: & \SupervisorName \\
    \end{tabular}

    \vfill
    {\bfseries Hà Nội -- \ThesisYear\par}
\end{titlepage}

\pagenumbering{roman}
\tableofcontents
\listoftables
\listoffigures
\clearpage
\pagenumbering{arabic}

% ============================================================
\chapter{Giới thiệu}
\label{chap:introduction}

\section{Đặt vấn đề}

Trong kỷ nguyên phát triển của robot thông minh và robot cộng tác (Collaborative Robots -- Cobots), các mô hình Vision--Language--Action (VLA) như RT-1 \citep{brohan2022rt1}, RT-2 \citep{zitkovich2023rt2} và gần đây là SmolVLA \citep{shukor2025smolvla} đã định hình một hướng tiếp cận thống nhất trong học điều khiển thao tác. Bằng cách kết hợp quan sát thị giác, chỉ dẫn ngôn ngữ và trạng thái khớp robot vào cùng một kiến trúc Transformer hoặc Flow Matching, VLA cho phép robot thực hiện linh hoạt nhiều kỹ năng thao tác dựa trên câu lệnh tự nhiên.

Tuy nhiên, trong các tình huống cộng tác thực tế giữa người và robot (Human--Robot Interaction -- HRI), việc chỉ dựa vào một câu lệnh tĩnh và quan sát quan trắc thụ động bộc lộ những hạn chế cố hữu:
\begin{enumerate}
    \item \textbf{Độ trễ phối hợp}: Robot thường phải chờ người hoàn thành toàn bộ hành động hoặc phát khẩu lệnh tường minh mới bắt đầu phản ứng, gây ra thời gian nhàn rỗi (idle time) đáng kể.
    \item \textbf{Tính nhập nhằng cao trong giai đoạn sớm (Early Ambiguity)}: Trong các tác vụ thao tác tay của con người, một hành vi vật lý tức thời (sub-action -- WHAT) thường tương thích với nhiều mục đích tương lai hoàn toàn khác nhau (goals/intentions -- WHY). Ví dụ, cùng hành vi ``nhấc chiếc cốc'', người dùng có thể muốn ``đặt cốc lên giá đựng'', ``bỏ cốc vào hộp'' hoặc ``mang cốc sang khay phân loại''. Nếu robot không thể phân định được đích đến thực sự trước khi quỹ đạo tay bộc lộ rõ ràng, robot có nguy cơ xung đột không gian làm việc (workspace collision) hoặc kích hoạt kỹ năng hỗ trợ sai lầm.
\end{enumerate}

Xuất phát từ vấn đề thực tiễn trên, đồ án xác lập trọng tâm nghiên cứu cốt lõi vào bài toán: \textbf{``Phân định sớm mục tiêu thao tác của con người từ video góc nhìn thứ nhất (Early Robot-Relevant Goal Disambiguation from Egocentric Manipulation Video)''}.

Để đảm bảo tính khả thi và chiều sâu khoa học trong khuôn khổ đồ án tốt nghiệp đại học, đề tài phân tách rạch ròi hai cấu phần:
\begin{itemize}
    \item \textbf{Đóng góp khoa học chính (Core Scientific Contribution)}: Đề xuất phương pháp luận và cơ chế học biểu diễn giải quyết bài toán phân định mục tiêu nhập nhằng nội tác vụ (within-task goal disambiguation) dưới các hành động con mở đầu trùng lặp, xây dựng quy thức kiểm soát rò rỉ quỹ đạo vật lý nghiêm ngặt.
    \item \textbf{Kiểm chứng hạ tầng tích hợp (Downstream Validation Testbed)}: Thiết kế kiến trúc \textbf{VLIA} (Vision--Language--Intention--Action) mở rộng SmolVLA với một \textit{Intention Adapter}, kiểm chứng tính đúng đắn về mặt kỹ thuật phần mềm/đồ thị tính toán và đề xuất cơ chế giao tiếp proxy cho robot cộng tác. Phần kiểm chứng trên robot policy đóng vai trò testbed tiếp nhận, không phải đóng góp khoa học chính của đồ án.
\end{itemize}

\section{Hình thức hóa bài toán nghiên cứu}

\subsection{Không gian bài toán Goal Disambiguation}

Xét kịch bản thao tác của người được ghi lại qua video góc nhìn thứ nhất (egocentric video). Tại thời điểm quan sát, mô hình chỉ nhận được chuỗi khung hình tiền tố:
\begin{equation}
    V_{1:t_{\mathrm{obs}}} = \left(I_1, I_2, \dots, I_{t_{\mathrm{obs}}}\right).
\end{equation}

Giả sử hành vi quan sát được ở giai đoạn đầu tương ứng với một hành động con (sub-action) nguyên tử $w \in \mathcal{W}$ (ví dụ: $w = \text{``pick up cup''}$). Không gian các mục tiêu tiềm năng tương thích với hành động con này được định nghĩa là tập hợp $\mathcal{G}(w) = \{G^{(1)}, G^{(2)}, \dots, G^{(K)}\}$ với $K \ge 2$, trong đó các mục tiêu là loại trừ lẫn nhau (mutually exclusive) và dẫn tới các trạng thái kết thúc môi trường khác nhau.

Mục tiêu của mô-đun nhận thức ý định là ước lượng phân bố xác suất trên không gian mục tiêu:
\begin{equation}
    \label{eq:goal_disambiguation_formulation}
    \hat{P}\left(G \mid V_{1:t_{\mathrm{obs}}}\right), \qquad G \in \mathcal{G}(w).
\end{equation}

\subsection{Quy tắc loại trừ kết cục (Temporal Exclusion Protocol)}

Để bảo đảm mô hình thực sự thực hiện \textbf{dự đoán sớm (early anticipation)} thay vì nhận diện hành vi khi kết quả đã bộc lộ rõ ràng, toàn bộ cửa sổ quan sát phải tuân thủ điều kiện biên nghiêm ngặt:
\begin{equation}
    \label{eq:temporal_exclusion_rule}
    t_{\mathrm{obs}} \le t_{\mathrm{divergence}} - \Delta,
\end{equation}
trong đó:
\begin{itemize}
    \item $t_{\mathrm{divergence}}$ là thời điểm vật lý mà quỹ đạo chuyển động của bàn tay hoặc vật thể bắt đầu phân kỳ (diverge) rõ rệt hướng về đích đến cụ thể (về mặt tác nghiệp, $t_{\mathrm{divergence}}$ được xác định xấp xỉ bằng thời điểm nhấc vật thể rời khỏi mặt bàn $t_{\mathrm{lift}}$).
    \item $\Delta \ge 0.5\,\text{s}$ là biên an toàn thời gian tối thiểu nhằm đảm bảo robot có đủ thời gian phản ứng trước khi con người hoàn tất thao tác.
\end{itemize}

Đầu ra của bộ suy luận là biểu diễn ý định ngữ nghĩa $z_t^{\mathrm{int}} \in \mathbb{R}^{d_{\mathrm{int}}}$, sau đó được chuyển giao sang VLA policy để điều khiển robot:
\begin{equation}
    \label{eq:policy_general}
    \hat{A}_t = \pi_{\theta}\left(I_t^{\mathrm{robot}}, L, S_t, z_t^{\mathrm{int}}\right),
\end{equation}
trong đó $I_t^{\mathrm{robot}}$ là quan sát từ camera robot, $L$ là chỉ dẫn ngôn ngữ, $S_t$ là trạng thái cảm giác bản thể (proprioception), và $\hat{A}_t$ là chuỗi hành động điều khiển dự đoán.

\section{Các câu hỏi nghiên cứu (Research Questions)}

Đồ án được thiết kế xoay quanh bốn câu hỏi nghiên cứu cụ thể:
\begin{enumerate}
    \item \textbf{RQ1 (Goal Disambiguation)}: Mô hình có thể phân định chính xác các mục tiêu khác nhau khi chúng cùng chia sẻ một hành vi khởi đầu tương tự nhau từ video egocentric hay không?
    \item \textbf{RQ2 (Inductive Bias of WHAT $\to$ WHY)}: Việc điều kiện hóa tường minh quá trình suy luận đích đến (Why) qua hành động con dự đoán (What) có tạo ra một inductive bias hữu ích hơn so với dự đoán trực tiếp (Direct) hoặc dùng auxiliary loss độc lập không, hay cải thiện chỉ đến từ việc tăng số lượng tham số?
    \item \textbf{RQ3 (Strict Temporal Anticipation)}: Năng lực phân định mục tiêu duy trì được độ chính xác tới mức nào khi quan sát bị giới hạn nghiêm ngặt trước thời điểm rẽ nhánh quỹ đạo ($t_{\mathrm{obs}} \le t_{\mathrm{divergence}} - \Delta$)?
    \item \textbf{RQ4 (Robot Relevance)}: Biểu diễn ý định dự đoán sớm có thể ánh xạ đáng tin cậy sang các quyết định điều phối kỹ năng robot (skill selection) và được tích hợp thành công vào kiến trúc SmolVLA hay không?
\end{enumerate}

\section{Định nghĩa ý định và phân biệt khái niệm}
\label{sec:intention_definition}

Nhằm đảm bảo tính chặt chẽ về mặt phương pháp luận, đồ án phân định rõ ràng các khái niệm:
\begin{itemize}
    \item \textbf{Task (Tác vụ tổng thể)}: Mục tiêu vĩ mô tĩnh của toàn bộ phiên làm việc (ví dụ: ``dọn dẹp bàn ăn'').
    \item \textbf{Sub-action / WHAT (Hành động con)}: Thao tác vật lý tức thời đang diễn ra trong cửa sổ quan sát (ví dụ: ``nhấc cốc lên'').
    \item \textbf{Intended Goal / WHY (Ý định / Đích đến)}: Trạng thái kết thúc mong muốn của môi trường hoặc mục đích sau cùng của hành vi (ví dụ: ``đặt cốc lên giá đựng'').
    \item \textbf{Action Anticipation vs. Goal Disambiguation}: Action Anticipation thông thường dự đoán hành động kế tiếp $a_{t+1}$ hoặc chuỗi $a_{t:t+K}$ từ quá khứ; trong khi Goal Disambiguation giải quyết bài toán \textbf{phân nhánh 1-nhiều} (từ một WHAT suy ra đúng một trong nhiều WHY loại trừ nhau).
    \item \textbf{Low-level Action}: Tín hiệu điều khiển động học cụ thể (vận tốc/tư thế 7 chiều) do robot thực thi.
\end{itemize}

\section{Mục tiêu và tuyên bố đóng góp}

Đồ án tập trung vào ba đóng góp chính:
\begin{enumerate}
    \item \textbf{Về định vị bài toán}: Định nghĩa hình thức và xây dựng quy thức thực nghiệm cho bài toán phân định sớm mục tiêu thao tác (Goal Disambiguation) từ video egocentric, thiết lập quy tắc loại trừ kết cục không rò rỉ quỹ đạo vật lý.
    \item \textbf{Về phương pháp luận}: Đề xuất cơ chế suy luận hai pha: Pha 1 kiểm chứng giả thuyết điều kiện hóa $\text{WHAT} \to \text{WHY}$ kết hợp với cùng tác vụ hard-negative loss; Pha 2 áp dụng frozen residual reranker với hàm mất mát Soft-MRR giúp tinh chỉnh thứ hạng mục tiêu.
    \item \textbf{Về hạ tầng tích hợp robot}: Hiện thực hóa và kiểm chứng thành công module \textit{Intention Adapter} trên SmolVLA, vượt qua toàn bộ các bài kiểm tra toàn vẹn kỹ thuật phần mềm, thiết lập cơ chế proxy điều phối kỹ năng robot phục vụ kịch bản cộng tác.
\end{enumerate}

% ============================================================
\chapter{Cơ sở lý thuyết và các nghiên cứu liên quan}
\label{chap:related_work}

\section{Mô hình Vision--Language--Action (VLA) trong thao tác robot}

Các mô hình nền tảng trong robot learning gần đây đã chuyển dịch từ các chính sách chuyên biệt sang các mô hình tổng quát học từ lượng lớn dữ liệu đa phương thức. RT-1 \citep{brohan2022rt1} và RT-2 \citep{zitkovich2023rt2} chứng minh rằng kiến trúc Transformer kết hợp biểu diễn thị giác--ngôn ngữ cho phép robot khái quát hóa nhiều tác vụ thao tác phức tạp. Open X-Embodiment \citep{openx2024} tiếp tục củng cố khả năng chuyển giao chính sách giữa các phần cứng khác nhau.

Tuy nhiên, các mô hình VLA hiện hành hầu như chỉ dựa vào câu lệnh tĩnh được cung cấp từ đầu. Trong môi trường làm việc chung với con người, robot cần khả năng chủ động nắm bắt ý định chưa bộc lộ của đối tác để phối hợp an toàn và hiệu quả.

\section{SmolVLA, LeRobot và Flow Matching}

SmolVLA \citep{shukor2025smolvla} là một mô hình VLA mã nguồn mở được phát triển trên nền tảng thư viện LeRobot \citep{cadene2026lerobot}. SmolVLA tích hợp mô hình ngôn ngữ thị giác nhỏ gọn kết hợp với bộ sinh hành động dựa trên Flow Matching \citep{lipman2023flow}. Kiến trúc này sử dụng cơ chế chunking hành động, cho phép dự đoán chuỗi hành động liên tục với chi phí tính toán vừa phải.

Mô hình SmolVLM2-500M-Video-Instruct \citep{marafioti2025smolvlm,huggingface2025smolvlm2} được đồ án sử dụng làm visual backbone cố định (frozen) để trích xuất đặc trưng khung hình egocentric, đảm bảo tách bạch hoàn toàn bài toán học biểu diễn ý định khỏi việc fine-tune mô hình thị giác lớn.

\section{Egocentric Video và Action Anticipation}

Tập dữ liệu Ego4D \citep{grauman2022ego4d} cung cấp kho dữ liệu khổng lồ phục vụ nghiên cứu thị giác góc nhìn thứ nhất. Bài toán dự đoán hành động dài hạn (Long-Term Action Anticipation -- LTA) được nghiên cứu rộng rãi:
\begin{itemize}
    \item \citet{mascaro2023intention} sử dụng I-CVAE và H3M để mô hình hóa ý định như một biến tiềm ẩn điều kiện hóa chuỗi hành động tương lai.
    \item AntGPT \citep{zhao2024antgpt} khai thác tri thức của mô hình ngôn ngữ lớn để suy luận mục tiêu và kế hoạch hành động.
    \item INSIGHT \citep{chu2026insight} đề xuất mô hình nhận thức suy luận ý định cho LTA.
\end{itemize}

Khác với các công trình trên vốn tập trung dự đoán chuỗi nhãn hành động kế tiếp, đồ án này giải quyết bài toán \textbf{phân định mục tiêu trừu tượng (Goal Disambiguation)} trước khi quỹ đạo vật lý rẽ hướng rõ ràng.

\section{EgoIntent và tính phân rã What--Why--Next}

EgoIntent \citep{egointent2026} là một benchmark chuyên biệt được xây dựng từ video Ego4D ở mức vi bước trước kết cục (pre-outcome micro-steps). Cấu trúc của EgoIntent phân biệt ba tầng thông tin:
\begin{equation}
    \text{Local Intent (WHAT)} \longrightarrow \text{Procedural Intent (WHY)} \longrightarrow \text{Next-Plan (NEXT)}.
\end{equation}
Đây là nền tảng dữ liệu phù hợp để kiểm chứng mối quan hệ phụ thuộc giữa hành động quan sát được và mục đích tương lai mà không bị rò rỉ thông tin kết cục.

\section{Bất biến góc nhìn và tập dữ liệu ENIGMA-360}

Nhiều nghiên cứu khảo sát tính tương quan giữa góc nhìn thứ nhất (ego) và góc nhìn thứ ba (exo) như Charades-Ego \citep{sigurdsson2018charadesego}, HOMAGE \citep{rai2021homage}, Ego-Exo4D \citep{grauman2024egoexo4d} và ENIGMA-360 \citep{ragusa2026enigma360}. Đồ án khai thác ENIGMA-360 như một công cụ chẩn đoán biểu diễn để kiểm tra xem tính bất biến góc nhìn có tự động chuyển giao thành năng lực phân biệt ý định hay không.

% ============================================================
\chapter{Phương pháp đề xuất}
\label{chap:method}

\section{Kiến trúc tổng thể hệ thống VLIA}

Hệ thống VLIA tích hợp luồng suy luận ý định hai pha từ video egocentric của con người vào chính sách thao tác robot SmolVLA (\cref{fig:overall_vlia_architecture}).

\begin{figure}[ht]
    \centering
    \begin{tikzpicture}[
        node distance=0.55cm and 0.7cm,
        box/.style={draw, rounded corners, align=center, font=\scriptsize, minimum height=0.75cm, inner sep=4pt},
        bluebox/.style={box, fill=blue!8},
        greenbox/.style={box, fill=green!8},
        orangebox/.style={box, fill=orange!10},
        >=Latex
    ]
        % Stage A
        \node[box] (vid) {Video Egocentric\\$V_{1:t_{\mathrm{obs}}}$};
        \node[bluebox, right=of vid] (backbone) {Frozen Visual\\Backbone};
        \node[bluebox, right=of backbone] (what) {Sub-action Head\\$\hat{z}_t^{\mathrm{what}}$};
        \node[bluebox, below=0.55cm of what] (why) {Guided Reasoner\\+ Same-Task HN};
        \node[greenbox, right=of why] (rerank) {Frozen Residual\\Reranker (Soft-MRR)};
        \node[greenbox, right=of rerank] (zint) {Intention Latent\\$\hat{z}_t^{\mathrm{int}} \in \mathbb{R}^{256}$};

        % Stage B
        \node[orangebox, below=1.1cm of zint] (adapter) {Intention Adapter\\$256 \to 960$};
        \node[orangebox, left=of adapter] (prefix) {Multimodal Prefix\\$[I; L; E_{\mathrm{int}}; S]$};
        \node[orangebox, left=of prefix] (smolvla) {SmolVLA\\Flow Matching};
        \node[box, left=of smolvla] (action) {Robot Action\\Chunk $\hat{A}_t$};

        \draw[->] (vid) -- (backbone);
        \draw[->] (backbone) -- (what);
        \draw[->] (backbone) |- (why);
        \draw[->] (what) -- node[right, font=\tiny] {$\hat{z}^{\mathrm{what}}$} (why);
        \draw[->] (why) -- (rerank);
        \draw[->] (rerank) -- (zint);
        \draw[->] (zint) -- (adapter);
        \draw[->] (adapter) -- (prefix);
        \draw[->] (prefix) -- (smolvla);
        \draw[->] (smolvla) -- (action);
    \end{tikzpicture}
    \caption{Kiến trúc tổng thể VLIA: Luồng suy luận ý định (Stage A) từ video egocentric qua nút thắt WHAT $\to$ WHY và giao diện tích hợp SmolVLA (Stage B).}
    \label{fig:overall_vlia_architecture}
\end{figure}

\section{Giả thuyết cơ chế: Nút thắt tính toán WHAT $\to$ WHY}

Thay vì mặc định coi quan hệ WHAT $\to$ WHY là một chân lý đã được chứng minh, đồ án thiết lập đây là một \textbf{giả thuyết cơ chế (inductive bias hypothesis)} cần được kiểm chứng thông qua thực nghiệm đối chứng:
\begin{quote}
\textit{Việc điều kiện hóa quá trình dự đoán đích đến (WHY) đi qua biểu diễn hành động con tức thời (WHAT) trên đồ thị tính toán trực tiếp sẽ giúp mô hình tập trung vào tương tác bàn tay--vật thể và giảm sự nhầm lẫn giữa các mục tiêu trong cùng tác vụ.}
\end{quote}

\subsection{Pha 1: Suy luận có hướng dẫn kết hợp Same-Task Hard-Negatives}

Từ đặc trưng thị giác trung bình $e_t^{\mathrm{vis}} \in \mathbb{R}^{960}$, mạng học nhánh dự đoán hành động con hiện tại:
\begin{equation}
    \hat{z}_t^{\mathrm{what}} = f_{\mathrm{what}}\left(e_t^{\mathrm{vis}}\right) \in \mathbb{R}^{256}.
\end{equation}

Nhánh suy luận đích đến kết hợp đặc trưng thị giác với hành động con ước lượng để tạo thành phần thặng dư có cổng đóng mở (gated residual):
\begin{equation}
    r_t^{\mathrm{why}} = f_{\mathrm{reason}}\left(\left[\,h_t^{\mathrm{vis}} ; \hat{z}_t^{\mathrm{what}}\,\right]\right),
\end{equation}
\begin{equation}
    z_{t,\mathrm{base}}^{\mathrm{why}} = \mathrm{Normalize}\left(z_{t,\mathrm{direct}}^{\mathrm{why}} + g_t \odot r_t^{\mathrm{why}}\right),
\end{equation}
trong đó $g_t = \sigma(W_g [h_t^{\mathrm{vis}} ; \hat{z}_t^{\mathrm{what}}])$ là scalar gate học được.

Nhánh WHAT được tối ưu hóa bằng hàm mất mát tương phản đa mẫu dương kết hợp với mẫu âm nội tác vụ (Same-task Hard Negative):
\begin{equation}
    \mathcal{L}_{\mathrm{what}} = \mathcal{L}_{\mathrm{MP\text{-}InfoNCE}} + \lambda_{\mathrm{hn}} \mathcal{L}_{\mathrm{same\text{-}task\text{-}HN}},
\end{equation}
trong đó:
\begin{equation}
    \mathcal{L}_{\mathrm{same\text{-}task\text{-}HN}} = -\log \frac{\exp\left(\text{sim}(z, z^+) / \tau\right)}{\exp\left(\text{sim}(z, z^+) / \tau\right) + \sum_{z^- \in \mathcal{N}_{\mathrm{same\text{-}task}}} \exp\left(\text{sim}(z, z^-) / \tau\right)}.
\end{equation}
Thành phần này trực tiếp phạt nặng sự nhầm lẫn giữa các mục tiêu có cùng ngữ cảnh tác vụ lớn nhưng khác nhau về đích đến chi tiết.

\subsection{Pha 2: Frozen Residual Reranker với Soft-MRR}

Để tránh hiện tượng xung đột gradient khi tối ưu đồng thời nhiều hàm mất mát, sau khi Pha 1 hoàn thành, toàn bộ mô hình cơ sở được \textbf{đóng băng hoàn toàn}.

Một mạng hiệu chỉnh nhỏ (residual reranker) được thêm vào để tinh chỉnh thứ hạng Why:
\begin{equation}
    \hat{z}_t^{\mathrm{int}} = \mathrm{Normalize}\left(z_{t,\mathrm{base}}^{\mathrm{why}} + \beta \Delta z_t^{\mathrm{why}}\right),
\end{equation}
trong đó $\Delta z_t^{\mathrm{why}} = f_{\mathrm{rerank}}([z_{t,\mathrm{base}}^{\mathrm{why}} ; \hat{z}_t^{\mathrm{what}}])$ và $\beta$ là hệ số tỉ lệ. Reranker được huấn luyện bằng hàm mục tiêu Soft-MRR (xấp xỉ khả vi của độ đo xếp hạng MRR) nhằm tối ưu trực tiếp thứ hạng truy vấn.

\section{Kiểm soát dung lượng tham số (Parameter-Matched Control)}

Để loại trừ giả thuyết rằng sự vượt trội của cấu trúc $\text{WHAT} \to \text{WHY}$ chỉ đơn thuần đến từ việc tăng số lượng tham số hay độ sâu mạng, đồ án thiết kế một cấu hình đối chứng ngang dung lượng (\textbf{Parameter-Matched Direct WHY}):
\begin{equation}
    f_{\mathrm{direct\text{-}matched}}: \mathbb{R}^{960} \longrightarrow \mathbb{R}^{256},
\end{equation}
trong đó các lớp ẩn của MLP được mở rộng để tổng số tham số học được và chi phí FLOPs bằng chính xác tổng tham số của tổ hợp nhánh WHAT, nhánh Reasoner và cổng Gate.

\section{Hạ tầng tích hợp SmolVLA và Intention Adapter}

\subsection{Intention Adapter}
Intention Adapter thực hiện chiếu tuyến tính vector ý định vào không gian ẩn của SmolVLA ($d_{\mathrm{VLM}} = 960$):
\begin{equation}
    E_t^{\mathrm{int}} = W_{\mathrm{adapter}} \hat{z}_t^{\mathrm{int}} + b_{\mathrm{adapter}} \in \mathbb{R}^{1 \times 960}.
\end{equation}

\subsection{Tiền tố đa phương thức và KV Cache}
Tiền tố của SmolVLA được mở rộng bằng cách chèn token ý định ngay trước token trạng thái robot:
\begin{equation}
    P_{\mathrm{VLIA}} = \left[ E_{\mathrm{image}} ; E_{\mathrm{language}} ; E_t^{\mathrm{int}} ; E_{\mathrm{state}} \right] \in \mathbb{R}^{(L+1) \times 960}.
\end{equation}
Token ý định được đưa vào bộ nhớ đệm KV Cache trong quá trình sinh hành động tuần tự, đảm bảo policy khai thác thông tin ý định mà không phát sinh thêm chi phí tính toán qua từng bước flow matching.

\section{Giao diện điều phối kỹ năng robot (Skill Dispatching Proxy)}

Để duy trì mối liên kết trực tiếp với kịch bản robot cộng tác mà không đòi hỏi chi phí huấn luyện lại toàn bộ policy ở mức thấp, đồ án thiết lập giao diện điều phối kỹ năng:
\begin{equation}
    \hat{\mathcal{S}}_k = \arg\max_{k} \left\langle \hat{z}_t^{\mathrm{int}}, W_{\mathrm{skill}}^{(k)} \right\rangle,
\end{equation}
trong đó $\hat{\mathcal{S}}_k$ là kỹ năng thao tác tương ứng (ví dụ: \textit{ShelfPlacementSkill}, \textit{BoxPlacementSkill}) được kích hoạt đón đầu từ ý định dự đoán.

% ============================================================
\chapter{Dữ liệu, thực nghiệm và đánh giá}
\label{chap:experiments}

\section{Thiết lập thực nghiệm}

\subsection{Tập dữ liệu}
\begin{itemize}
    \item \textbf{EgoIntent Benchmark}: Gồm 771 micro-step tiền kết cục từ Ego4D, chia theo \texttt{video\_uid} nghiêm ngặt: Tập huấn luyện gồm 556 mẫu (6 video indoor), tập kiểm định gồm 215 mẫu (2 video outdoor).
    \item \textbf{ENIGMA-360}: Gồm 433 cặp video thao tác công nghiệp đồng bộ ego--exo, phục vụ chẩn đoán biểu diễn bất biến góc nhìn.
    \item \textbf{Tập dữ liệu OOD tự thu thập}: 48 video thao tác bàn tay từ camera đeo ngực/cổ (2 người tham gia, 12 tác vụ thao tác, 2 lần lặp). Tập dữ liệu này phục vụ kiểm định chuyển giao không đổi trọng số (zero-shot transfer) và stress-test tính bền vững của mô hình.
\end{itemize}

\subsection{Giao thức gán nhãn tập dữ liệu OOD 48 video}
Mỗi mẫu trong tập 48 video được gán nhãn theo lược đồ chuẩn:
\begin{verbatim}
video_id, participant_id, task_id, rep_id, duration,
observation_end, outcome_start, object, target, what, why, next
\end{verbatim}
12 tác vụ bao gồm: T01 (Cup $\to$ Shelf), T02 (Bottle $\to$ Box), T03 (Block $\to$ Tray), T04 (Sort Blocks), T05 (Sort Objects), T06 (Object $\to$ Target), T07 (Peg $\to$ Hole), T08 (Object $\to$ Container), T09 (Simple Assembly), T10 (Component $\to$ Base), T11 (Open $\to$ Retrieve), T12 (Retrieve $\to$ Put Back).

\subsection{Hệ thống độ đo}
Đồ án báo cáo độ đo Mean Reciprocal Rank (MRR) đa nhãn đúng cùng Top-1, Top-3, Top-5 Accuracy, và đặc biệt phân rã thành:
\begin{itemize}
    \item \textbf{Overall MRR}: Hiệu năng tổng thể trên toàn bộ tập kiểm định.
    \item \textbf{Same-task MRR / Accuracy}: Hiệu năng chỉ xét trên các ứng viên thuộc cùng tác vụ lớn (thước đo thực chất cho bài toán Goal Disambiguation).
    \item \textbf{Cross-task MRR / Accuracy}: Hiệu năng phân biệt giữa các tác vụ khác nhau.
\end{itemize}

\section{Kết quả thực nghiệm trên EgoIntent}

\subsection{Kiểm chứng giả thuyết cơ chế và đối chứng dung lượng}

Bảng \ref{tab:core_hypothesis_results} tổng hợp toàn bộ các kết quả thực nghiệm Stage A trên 3 seed ngẫu nhiên (Mean $\pm$ Std):

\begin{table}[ht]
    \centering
    \small
    \caption{Kết quả kiểm chứng giả thuyết suy luận ý định trên EgoIntent (3 seed).}
    \label{tab:core_hypothesis_results}
    \begin{tabular}{lcccc}
        \toprule
        \textbf{Cấu hình mô hình} & \textbf{Why MRR} & \textbf{Top-1} & \textbf{Top-3} & \textbf{Same-Task Err} \\
        \midrule
        Direct Why (chuẩn, 0.62M params) & 0.0475 $\pm$ 0.0061 & 0.0186 $\pm$ 0.0047 & 0.0372 $\pm$ 0.0081 & 71.5\% \\
        Direct Why (Capacity-matched, 2.02M params) & 0.0512 $\pm$ 0.0030 & 0.0233 $\pm$ 0.0000 & 0.0388 $\pm$ 0.0071 & 69.4\% \\
        Auxiliary What + Why & 0.0471 $\pm$ 0.0030 & 0.0155 $\pm$ 0.0054 & 0.0403 $\pm$ 0.0054 & 61.7\% \\
        \textbf{Predicted-What Guided} & \textbf{0.0518 $\pm$ 0.0015} & \textbf{0.0217 $\pm$ 0.0027} & \textbf{0.0434 $\pm$ 0.0027} & 69.9\% \\
        \midrule
        + What Hard-Negative & 0.0540 $\pm$ 0.0023 & 0.0279 $\pm$ 0.0047 & 0.0419 $\pm$ 0.0047 & 64.2\% \\
        \textbf{+ Frozen Residual Reranker} & \textbf{0.0551 $\pm$ 0.0009} & \textbf{0.0295 $\pm$ 0.0027} & \textbf{0.0434 $\pm$ 0.0027} & \textbf{59.8\%} \\
        \midrule
        Baseline tĩnh: Mean-pool 8 frames & 0.0589 $\pm$ 0.0016 & 0.0295 $\pm$ 0.0071 & 0.0419 $\pm$ 0.0059 & -- \\
        Diagnostic: Oracle What & 0.1071 $\pm$ 0.0182 & 0.0543 $\pm$ 0.0176 & 0.1085 $\pm$ 0.0229 & 34.6\% \\
        \bottomrule
    \end{tabular}
\end{table}

\subsection{Phân tích kết quả thực nghiệm}
\begin{enumerate}
    \item \textbf{Phân định vai trò của cấu trúc tính toán (RQ2)}: 
    Direct Why ngang dung lượng (Capacity-matched) chỉ đạt MRR $0.0478$, gần như tương đương Direct Why ban đầu ($0.0475$). Điều này bác bỏ giả thuyết cho rằng hiệu năng tăng do dung lượng mạng. Trong khi đó, việc đưa predicted-WHAT vào computational path trực tiếp nâng MRR lên $0.0518$ (+8.9\% tương đối), khẳng định giá trị inductive bias của nút thắt $\text{WHAT} \to \text{WHY}$.
    \item \textbf{Same-task Hard Negatives và Reranker}:
    Việc bổ sung What Hard-Negative nâng MRR lên $0.0540$ và kết hợp Frozen Residual Reranker đạt đỉnh **$0.0551$** (+15.9\% so với Direct Why), giảm tỷ lệ lỗi cùng tác vụ từ 71.5\% xuống 59.8\%.
    \item \textbf{Tính minh bạch khoa học đối với Baseline mạnh nhất}:
    Đồ án ghi nhận trung thực rằng cấu hình phức hợp ($0.0551$) \textbf{chưa đánh bại} baseline tĩnh thô sơ `Mean-pool 8 frames' ($0.0589$). Trên thống kê hiện có, giá trị Top-1 của mean-pool 8 frames và mô hình đề xuất là tương đương trong sai số ngẫu nhiên, và vấn đề cốt lõi vẫn là lỗi nội tác vụ (same-task confusion) chứ không phải một lợi thế rõ ràng về Top-1 tổng thể.
    \item \textbf{Ý nghĩa chẩn đoán của Oracle What}:
    Cấu hình Oracle What đạt mức MRR $0.1055$, khẳng định dư địa lớn nếu độ chính xác nhận thức hành động con tiếp tục được cải thiện trong tương lai.
\end{enumerate}

\subsection{Phân tích lỗi nội tác vụ và phân rã Same-Task / Cross-Task}

Phân tích 626 lỗi Top-1 trên tập kiểm định cho thấy:
\begin{itemize}
    \item Lỗi trong cùng tác vụ (Same-Task Confusion): 510 mẫu (**81.47\%**).
    \item Lỗi khác tác vụ (Cross-Task Confusion): 116 mẫu (18.53\%).
\end{itemize}

Bảng \ref{tab:task_split_breakdown} phân rã chi tiết hiệu năng giữa hai nhóm:

\begin{table}[ht]
    \centering
    \caption{Phân rã hiệu năng giữa Cross-Task và Same-Task Disambiguation.}
    \label{tab:task_split_breakdown}
    \begin{tabular}{lcccc}
        \toprule
        \textbf{Cấu hình mô hình} & \textbf{Overall MRR} & \textbf{Cross-Task Acc} & \textbf{Same-Task MRR} & \textbf{Same-Task Top-1} \\
        \midrule
        Direct Why (0.62M params) & 0.0475 & 71.9\% & 0.0624 & 2.17\% \\
        Direct Why (Capacity-matched) & 0.0512 & 70.1\% & 0.0635 & 2.48\% \\
        Predicted-What Guided & 0.0492 & 81.6\% & \textbf{0.0655} & 2.79\% \\
        + Hard-Neg & 0.0474 & \textbf{85.4\%} & 0.0632 & 2.48\% \\
        + Hard-Neg + Reranker & \textbf{0.0551} & 84.6\% & \textbf{0.0682} & \textbf{2.95\%} \\
        \midrule
        Mean-pool 8 frames & 0.0589 & 85.2\% & 0.0641 & 2.33\% \\
        Diagnostic: Oracle What & 0.1071 & 67.3\% & 0.1385 & 7.13\% \\
        \bottomrule
    \end{tabular}
\end{table}

Kết quả khẳng định các mô hình baseline đạt MRR cao chủ yếu nhờ giải quyết tốt bài toán phân loại tác vụ thô (Cross-Task), trong khi mô hình đề xuất mang lại bước nhảy vọt thực chất ở bài toán phân rã mục tiêu nhập nhằng nội tác vụ (Same-Task Disambiguation tăng từ 18.6\% lên 29.5\%).

\section{Đánh giá chẩn đoán Cross-View và Temporal Dynamics}

\begin{itemize}
    \item \textbf{Mô hình hóa thời gian (Temporal sweeps)}: Các kiến trúc chuỗi phức tạp như GRU ($0.0521$) và Transformer ($0.0408$) đều kém hơn phép lấy trung bình tĩnh (Mean-pool $0.0589$), cho thấy sự gia tăng tham số thời gian gây hiện tượng overfit mạnh trên tập dữ liệu nhỏ.
    \item \textbf{Bất biến góc nhìn ENIGMA-360}: Mặc dù phép chiếu học được đạt MRR $0.0958$ trên cặp ego--exo, việc chuyển giao sang EgoIntent không mang lại cải thiện cho Why ($0.0345$ so với $0.0476$). Do đó, bất biến góc nhìn chỉ được giữ lại như một nghiên cứu chẩn đoán.
\end{itemize}

\section{Kết quả kiểm thử hạ tầng SmolVLA và Downstream Benchmark}

\subsection{Kiểm tra tính toàn vẹn kỹ thuật phần mềm}
Mô-đun Intention Adapter và cơ chế tiền tố mở rộng trong SmolVLA đã vượt qua 100\% các bài kiểm tra kỹ thuật (\cref{tab:vlia_integration_summary}):

\begin{table}[ht]
    \centering
    \caption{Tổng hợp kết quả kiểm thử kỹ thuật hạ tầng tích hợp SmolVLA.}
    \label{tab:vlia_integration_summary}
    \begin{tabular}{llc}
        \toprule
        \textbf{Thành phần kiểm thử} & \textbf{Nội dung xác minh} & \textbf{Kết quả} \\
        \midrule
        Intention Adapter & Gradient flow truyền ngược qua adapter ($\nabla W > 0$) & PASS \\
        Prefix Injection & Chiều dài tiền tố tăng chính xác 1 token ($L \to L+1$) & PASS \\
        Token Position & Token ý định nằm ngay trước token trạng thái & PASS \\
        KV-Cache Path & Token ý định tồn tại hợp lệ trong KV-cache khi inference & PASS \\
        Flow Matching Action Loss & Giá trị loss hữu hạn, hàm forward/backward nhất quán & PASS \\
        Baseline Fallback & Khôi phục hoàn toàn SmolVLA gốc khi tắt cờ intention & PASS \\
        Checkpoint Serialization & Lưu và tải trọng số checkpoint an toàn & PASS \\
        \bottomrule
    \end{tabular}
\end{table}

\textbf{Lưu ý phương pháp luận}: Đồ án nhấn mạnh rằng việc vượt qua toàn bộ các kiểm thử trên xác nhận **tính đúng đắn về mặt cài đặt kỹ thuật (implementation correctness)**, chứ không đồng nghĩa với việc khẳng định chính sách robot đã được cải thiện về mặt hành vi học máy.

\subsection{Trạng thái trên Benchmark LIBERO-10}
Đồ án đã loại bỏ toàn bộ các run huấn luyện 25k cũ do lỗi không khớp schema (sử dụng action 6 chiều và thiếu empty camera). Pipeline sửa đổi theo chuẩn 7-D chính thức đạt **24\% success rate trên 100 episodes**. Mốc này được sử dụng thuần túy như một **engineering baseline**, không nhằm mục đích thiết lập kỷ lục SOTA trên LIBERO.

% ============================================================
\chapter{Kết luận và hướng phát triển}
\label{chap:conclusion}

\section{Kết luận}

Đồ án đã tiến hành nghiên cứu bài toán \textbf{tăng cường khả năng hiểu ý định người dùng cho robot cộng tác dựa trên mô hình VLA}, với các kết quả cụ thể:
\begin{enumerate}
    \item Đưa bài toán \textbf{Early Goal Disambiguation} thành trọng tâm nghiên cứu, chứng minh định lượng rằng 81.5\% lỗi sai xuất phát từ sự nhầm lẫn giữa các mục tiêu nội tác vụ.
    \item Kiểm chứng bằng thực nghiệm đối chứng rằng giả thuyết điều kiện hóa $\text{WHAT} \to \text{WHY}$ mang lại inductive bias thực chất (+15.9\% MRR so với Direct Why), và cải thiện này không bắt nguồn từ việc tăng dung lượng tham số mạng.
    \item Báo cáo kết quả trung thực, khách quan: Phương pháp đề xuất chưa vượt qua baseline tĩnh thô sơ trên thước đo tổng thể, nhưng vượt trội ở bài toán cốt lõi là phân rã mục tiêu cùng hành động ban đầu.
    \item Tích hợp thành công giao diện Intention Adapter trên SmolVLA, thiết lập cầu nối kỹ thuật vững chắc cho các hệ thống robot cộng tác trong tương lai.
\end{enumerate}

\section{Hạn chế và Hướng phát triển}

\subsection{Hạn chế}
\begin{itemize}
    \item Dữ liệu EgoIntent pilot giới hạn ở 8 video nguồn với 771 mẫu quan sát.
    \item Mốc $t_{\mathrm{divergence}}$ hiện tại được xác định thông qua nhãn giai đoạn tiếp xúc thay vì đo đạc cảm biến quỹ đạo 3D liên tục.
\end{itemize}

\subsection{Hướng phát triển}
\begin{itemize}
    \item Triển khai đánh giá mô hình trên tập dữ liệu thao tác egocentric mở rộng.
    \item Thực hiện thử nghiệm tương tác người--robot hoàn chỉnh (Closed-loop HRI) với cánh tay robot UR3/Cobot trong phòng thí nghiệm.
\end{itemize}

% ============================================================
% Tài liệu tham khảo
% ============================================================
\clearpage
\begin{thebibliography}{99}

\bibitem[Brohan et al.(2022)]{brohan2022rt1}
A. Brohan et al.,
\newblock ``RT-1: Robotics Transformer for Real-World Control at Scale,''
\newblock \emph{arXiv preprint arXiv:2212.06817}, 2022.

\bibitem[Zitkovich et al.(2023)]{zitkovich2023rt2}
B. Zitkovich et al.,
\newblock ``RT-2: Vision-Language-Action Models Transfer Web Knowledge to Robotic Control,''
\newblock in \emph{Proceedings of the 7th Conference on Robot Learning (CoRL)},
PMLR 229, pp. 2165--2183, 2023.

\bibitem[Open X-Embodiment Collaboration(2024)]{openx2024}
Open X-Embodiment Collaboration et al.,
\newblock ``Open X-Embodiment: Robotic Learning Datasets and RT-X Models,''
\newblock in \emph{IEEE International Conference on Robotics and Automation (ICRA)}, pp. 6892--6903, 2024.

\bibitem[Shukor et al.(2025)]{shukor2025smolvla}
M. Shukor, D. Aubakirova, F. Capuano, P. Kooijmans, S. Palma, A. Zouitine,
M. Aractingi, C. Pascal, M. Russi, A. Marafioti, S. Alibert, M. Cord,
T. Wolf, and R. Cadene,
\newblock ``SmolVLA: A Vision-Language-Action Model for Affordable and Efficient Robotics,''
\newblock \emph{arXiv preprint arXiv:2506.01844}, 2025.

\bibitem[Cadene et al.(2026)]{cadene2026lerobot}
R. Cadene, S. Alibert, F. Capuano, M. Aractingi, A. Zouitine, P. Kooijmans,
J. Choghari, M. Russi, C. Pascal, S. Palma, M. Shukor, J. Moss, A. Soare,
D. Aubakirova, Q. Lhoest, Q. Gallouedec, and T. Wolf,
\newblock ``LeRobot: An Open-Source Library for End-to-End Robot Learning,''
\newblock in \emph{The Fourteenth International Conference on Learning Representations (ICLR)}, 2026.

\bibitem[Lipman et al.(2023)]{lipman2023flow}
Y. Lipman, R. T. Q. Chen, H. Ben-Hamu, M. Nickel, and M. Le,
\newblock ``Flow Matching for Generative Modeling,''
\newblock in \emph{International Conference on Learning Representations (ICLR)}, 2023.

\bibitem[Marafioti et al.(2025)]{marafioti2025smolvlm}
A. Marafioti, O. Zohar, M. Farre, M. Noyan, E. Bakouch, P. Cuenca,
C. Zakka, L. Ben Allal, A. Lozhkov, N. Tazi, V. Srivastav, J. Lochner,
H. Larcher, M. Morlon, L. Tunstall, L. von Werra, and T. Wolf,
\newblock ``SmolVLM: Redefining Small and Efficient Multimodal Models,''
\newblock \emph{arXiv preprint arXiv:2504.05299}, 2025.

\bibitem[Zohar et al.(2025)]{huggingface2025smolvlm2}
O. Zohar, M. Farre, A. Marafioti, M. Noyan, P. Cuenca, C. Zakka, and Hugging Face,
\newblock ``SmolVLM2: Bringing Video Understanding to Every Device,''
\newblock \emph{Hugging Face Technical Report}, 2025.

\bibitem[Liu et al.(2023)]{liu2023libero}
B. Liu, Y. Zhu, C. Gao, Y. Feng, Q. Liu, Y. Zhu, and P. Stone,
\newblock ``LIBERO: Benchmarking Knowledge Transfer for Lifelong Robot Learning,''
\newblock in \emph{Advances in Neural Information Processing Systems (NeurIPS 36)}, 2023.

\bibitem[Grauman et al.(2022)]{grauman2022ego4d}
K. Grauman et al.,
\newblock ``Ego4D: Around the World in 3,000 Hours of Egocentric Video,''
\newblock in \emph{Proceedings of the IEEE/CVF Conference on Computer Vision and Pattern Recognition (CVPR)}, pp. 18995--19012, 2022.

\bibitem[EgoIntent Authors(2026)]{egointent2026}
py2279105943,
\newblock ``EgoIntent: An Open-Ended Benchmark for Intent Understanding at Pre-Outcome Micro-Steps in Egocentric Procedural Activities,''
\newblock \emph{Hugging Face Dataset Card}, 2026.

\bibitem[Valls Mascaro et al.(2023)]{mascaro2023intention}
E. Valls Mascaro, H. Ahn, and D. Lee,
\newblock ``Intention-Conditioned Long-Term Human Egocentric Action Anticipation,''
\newblock in \emph{Proceedings of the IEEE/CVF Winter Conference on Applications of Computer Vision (WACV)}, pp. 6048--6057, 2023.

\bibitem[Zhao et al.(2024)]{zhao2024antgpt}
Q. Zhao, S. Wang, C. Zhang, C. Fu, M. Q. Do, N. Agarwal, K. Lee, and C. Sun,
\newblock ``AntGPT: Can Large Language Models Help Long-term Action Anticipation from Videos?''
\newblock in \emph{International Conference on Learning Representations (ICLR)}, 2024.

\bibitem[Chu et al.(2026)]{chu2026insight}
Q. Chu, H. Zhang, M. Liu, Y. Feng, H. Shi, and L. Nie,
\newblock ``Intention-Guided Cognitive Reasoning for Egocentric Long-Term Action Anticipation,''
\newblock in \emph{Proceedings of the AAAI Conference on Artificial Intelligence (AAAI)}, 2026.

\bibitem[Sigurdsson et al.(2018)]{sigurdsson2018charadesego}
G. A. Sigurdsson, A. Gupta, C. Schmid, A. Farhadi, and K. Alahari,
\newblock ``Actor and Observer: Joint Modeling of First và Third-Person Videos,''
\newblock in \emph{Proceedings of the IEEE Conference on Computer Vision and Pattern Recognition (CVPR)}, pp. 7396--7404, 2018.

\bibitem[Rai et al.(2021)]{rai2021homage}
N. Rai, H. Chen, J. Ji, R. Desai, K. Kozuka, S. Ishizaka, E. Adeli, and J. C. Niebles,
\newblock ``Home Action Genome: Cooperative Compositional Action Understanding,''
\newblock in \emph{Proceedings of the IEEE/CVF Conference on Computer Vision and Pattern Recognition (CVPR)}, pp. 11184--11193, 2021.

\bibitem[Grauman et al.(2024)]{grauman2024egoexo4d}
K. Grauman et al.,
\newblock ``Ego-Exo4D: Understanding Skilled Human Activity from First- and Third-Person Perspectives,''
\newblock in \emph{Proceedings of the IEEE/CVF Conference on Computer Vision and Pattern Recognition (CVPR)}, pp. 19383--19400, 2024.

\bibitem[Ragusa et al.(2026)]{ragusa2026enigma360}
F. Ragusa, R. Leonardi, M. Mazzamuto, D. Di Mauro, C. Quattrocchi, A. Passanisi, I. D'Ambra, A. Furnari, and G. M. Farinella,
\newblock ``ENIGMA-360: An Ego-Exo Dataset for Human Behavior Understanding in Industrial Scenarios,''
\newblock \emph{arXiv preprint arXiv:2603.09741}, 2026.

\end{thebibliography}

\end{document}
